\section{Class Dichotomy, Characterization, and Escapes}

Fix an admissible parameter $a$, put
\[
 A=1+4a^2,\qquad \delta_\tau=1-A\tau^2,
 \qquad \alpha=-\delta_\tau A,
\]
and first consider the prime-member family $b=\varepsilon q$, with
$\varepsilon\in\{\pm1\}$ and $q$ prime.  An \emph{aligned class} means a
nonempty prime progression in which the Hilbert symbols at every controlled
finite place, at the moving place $q$, and at the real place are all $+1$.
This is deliberately a frozen-side notion: odd-valuation primes emerging from
the value of $P(b)$ are a separate obligation.
% src: THEOREMS.md:973-991

\subsection{The L10 dichotomy}

\begin{proposition}[Canonical alignment and the opposite wall; \textnormal{PROVED in the stated scope}]
On the canonical branch $\tau=2a/A$, the prime-$b$ family is aligned with the
canonical square class $-1$ in the exact sense that
\[
 \delta_\tau=\frac1A,\qquad \alpha=-1.
\]
On the branch $\tau=0$, if $a$ is odd, $A$ is prime,
$A\nmid 2zD_z$, and the prime $q\ne A$ satisfies the prescribed unit
conditions, then no aligned class exists.
\end{proposition}
% src: THEOREMS.md:1003-1014
% src: THEOREMS.md:1016-1042

The first assertion is an identity, not permission to discard the places
above $A$: although $\alpha$ is an $A$-unit, $\delta_\tau=1/A$ forces those
places back into the controlled set.  This distinction is essential; the
obsolete grid count obtained by omitting them is not used here.  For the
second assertion, $\tau=0$ gives $\delta_\tau=1$ and $\alpha=-A$.  The local
calculation at $A$ yields
\[
 (x,d)_A\,\leg{A}{q}=\leg{-2\varepsilon}{A}=-1.
\]
Indeed admissibility gives $A\equiv5\pmod 8$, so
$\leg{-1}{A}=+1$ and $\leg{2}{A}=-1$, independently of the sign
$\varepsilon$.  Two required symbols are therefore permanently
anti-correlated.  This is a theorem for prime integral $A$ under the displayed
hypotheses; it is not a claim extracted from a finite failure search.
% src: THEOREMS.md:1016-1042

\subsection{Free branches, an immovable wall}

Branch completion itself costs nothing.  If $\tau=p/q$ is rational in the
existing parameters and its denominator is proved nonzero on the base
locus, adjoining that branch introduces no inverse witness.  Admissibility
makes $A$ a nonsquare in $\Q_2$, hence $\delta_\tau\ne0$; the norm identity
in the unrefereed argument of Sun~\cite{sun} is uniform in $\tau$; and a
finite disjunction of branch equations is represented by their product in
the same witnesses.  Explicit choices meet all four square classes of
$\langle-1,A\rangle$, including the square branch
\[
 \tau^\dagger=\frac{1+2a^2}{1+4a^2},\qquad
 \alpha=(2a^2)^2.
\]
% src: THEOREMS.md:1100-1146
% src: RESULTS.md:89-91

The corrected controlled set for a prime member is
\[
 S=\{2,3,5,7\}\cup\supp(\alpha)\cup\supp(\delta_\tau).
\]
Outside $S$, the fixed content of $P$ has even valuation; after that content
is removed, the reduction is a nonzero polynomial whose degree prevents a
permanent odd valuation at the remaining primes.  Thus no fixed bad place is
silently omitted.
% src: THEOREMS.md:1189-1209

More importantly, changing branches does not change the obstruction.  For
prime $A$, admissible $a$, and $v_A(z)\le0$, every branch in the relevant
square-class group satisfies
\[
 (x,d)_A(x,d)_q=-1.
\]
Hence at least one controlled symbol is negative.  This branch-independent
wall was checked exactly on $2{,}136$ extended instances across all four
explicit branches; the computation confirms the proved identity rather than
replacing it.
% src: THEOREMS.md:1214-1239
% src: RESULTS.md:95

\begin{theorem}[The L11h characterization; \textnormal{PROVED}]
For the prime-$b$ family, an aligned class exists at the cell $(w,u)$, with
$z=wu$, if and only if the numerator of $z$ has a prime divisor
$p\equiv1\pmod4$.  On the finite grid this selects $190$ of $353$ cells; the
other $163$ cells are unreachable by this family on every branch.
\end{theorem}
% src: THEOREMS.md:1241-1265
% src: RESULTS.md:96

\begin{proof}[Proof sketch]
For necessity, the generalized wall says that alignment is impossible unless
some odd prime $p\mid A$ also has $v_p(z)>0$.  Write $a=m/n$ in lowest terms.
Such a prime divides $n^2+4m^2$, and therefore
$(2m n^{-1})^2\equiv-1\pmod p$; consequently $p\equiv1\pmod4$.  The
factor-by-factor telescoping argument applies also when $A$ is composite or
rational, so this necessity is \textsc{proved} for every admissible $a$, not
merely supported by a search.
% src: THEOREMS.md:1415-1431

Conversely, let $p\equiv1\pmod4$ divide the numerator of $z$.  Choose odd
representatives $m,n$ satisfying
$n^2\equiv-4m^2\pmod p$ and set $a=m/n$.  Then $a$ is admissible and
$v_p(A)>0$, so the factor shared by $A$ and $z$ breaks the wall.  Freezing the
complete controlled set gives an explicit aligned progression.  The
construction, including its escape prime, chosen $a$, branch, and modulus,
is replayed on all $190$ predicted rows.  This constructive sufficiency and
the preceding necessity make the equivalence the structural centre of the
class analysis.
% src: THEOREMS.md:1253-1259
\end{proof}

\subsection{Non-coprime escapes and complete grid coverage}

The characterization concerns $b=\varepsilon q$.  The only effective door
through the generalized coprime wall is to share a cell prime deliberately:
\[
 b=\varepsilon f q,\qquad f\mid z.
\]
This produces aligned escape classes at $103$ of the $163$ prime-family-walled
cells.  Their frozen-side alignment is \textsc{proved} row by row; the
remaining $60$ walled cells with no such escape are closed instead by direct,
fully soluble witnesses.
% src: RESULTS.md:98-100

The Schinzel audit also keeps six legacy residual rational/composite-$b$
witnesses separate as pointwise certificates.  They are unconditional local
witnesses, but they are not promoted into prime progressions and give no
prime-value theorem.  The eventual last cell, $(89,-2)$, was closed by the
replayed square-branch witness
$(a,b,\tau)=(3,89/367,19/37)$, completing cell coverage.
% src: README.md:40
% src: THEOREMS.md:1374-1382

\begin{center}
\begin{tabular}{lcc}
\hline
verified object & coverage & status\\
\hline
L11 prime-$b$ aligned classes & $190/190$ & \textsc{proved per row}\\ % src: RESULTS.md:96
non-coprime escape classes & $103/103$ & \textsc{proved per row}\\ % src: RESULTS.md:98-99
all aligned classes, with a zero-bad member & $293/293$ & \textsc{proved per row}\\ % src: RESULTS.md:108
all grid cells, with a soluble witness & $353/353$ & \textsc{proved per row}\\ % src: RESULTS.md:100
\hline
\end{tabular}
\end{center}
The class total is exactly the disjoint sum of the L11 and escape families;
the cell total additionally uses direct pointwise witnesses.  Neither total
asserts a uniform theorem beyond the verified grid.
% src: THEOREMS.md:1481-1509

\subsection{Why no square-class shortcut remains}

\begin{proposition}[L13c; \textnormal{PROVED per row}]
For every one of the $706$ cell--branch rows, $P$ is irreducible of degree
$8$ over $\Q$.  Consequently no constant-square or weak square-factor form
collapses the emergent-place obligation on the grid.
\end{proposition}
% src: THEOREMS.md:1384-1393
% src: RESULTS.md:101

The boundary between theorem and experiment is sharp.  The irreducibility
certificates are \textsc{proved}.  By contrast, the $137.5$ million exact
square-class tests with no admissible hit are \textsc{evidence} against a
hidden special-value shortcut, not a universal nonexistence theorem.
Likewise the observed Galois group $D_8\wr C_2$ is \textsc{evidence}: its
cycle types matched all $28{,}240$ certified unramified Frobenius samples,
but no group-identification proof is claimed.
% src: THEOREMS.md:1393-1402
% src: RESULTS.md:101

Finally, the earlier zero-bad search already supplied $24$ certified members
across $15$ distinct cells; those are \textsc{proved instances}, while the
rough quarter-density inferred from the sample is only \textsc{evidence}.
The later per-class replay strengthens the finite conclusion to the
$293/293$ line in the table, without changing the open uniform problem.
% src: THEOREMS.md:1404-1413
% src: THEOREMS.md:1481-1510
